\capitulo{6}{Trabajos relacionados}

En este apartado se exponen los artículos científicos relacionados con el TFG que han sido publicados hasta la fecha. 
\section{TFG relacionados}
\subsection{Retrieval-Augmented Generation para la Extracción de Información de Documentos Inteligente}
En este TFG se define y materializa un sistema RAG aplicado al ámbito clínico para extender las capacidades de los LLM existentes. Como arquitectura propone una canalización de recuperación y generación que indexa documentación médica, crea embeddings y orquesta prompts para aportar contexto fiable al modelo generativo. Incluye mejoras de retrieval mediante chunking con overlap y un MultiQuery Retriever para incrementar cobertura y precisión antes de la síntesis con un LLM. Implementa la solución sobre una API de Azure OpenAI, LangChain con FAISS como índice vectorial, y una interfaz web en Flask que simula el flujo real de consulta del paciente. Para su evaluación emplea RAGAS para medir la precisión, relevancia y coherencia y así determinar la seleccionar final~\cite{trabajorelacionado:Daniel}. 